\documentclass[11pt]{article}
\usepackage[margin=1in]{geometry}
\usepackage[T1]{fontenc}
\usepackage[utf8]{inputenc}
\usepackage{hyperref}
\hypersetup{colorlinks=true,urlcolor=blue}
\title{English Summaries of Working Papers\\{\large Ke Cheng Hsu}}
\author{Musashi University faculty application (D126070158)\\
\textit{Per the posting, existing abstracts need not be rewritten to 400 words.}}
\date{}
\begin{document}
\maketitle
\vspace{-1em}

\section*{Working papers}

\subsection*{1. Designing Return Policy in E-Commerce: A Trade-off between Purchase and Return Efficiency (JMP)}
We study whether return policies encourage efficient returns and purchases. A stricter return policy curbs inefficient ex-post returns. The resulting lower return rates, however, make higher prices more attractive for sellers and can exclude ex-ante efficient purchases. We compare automatic coverage that reimburses return-shipping costs (RI) with optional coverage (CI). We estimate a structural model of purchase, insurance, and return decisions under monopoly pricing across six apparel categories. A lenient policy such as RI performs better in high-margin, high-salvage categories by expanding valuable purchases. A strict policy such as CI performs better in thinner-margin categories by curbing excessive returns. Partial return-penalty discount (RPD)---a menu that pairs a discounted price with a medium restocking fee---performs better than RI and CI for most categories by balancing purchase and return efficiency more effectively.

\subsection*{2. Product Design with Heterogeneous Search Costs (with Pak Hung Au)}
Building on Wolinsky's (1986) sequential search model with differentiated products, this paper introduces search-cost heterogeneity. Firms use observable product designs to influence consumer search. In equilibrium, search patterns are partitioned: consumers with low search costs tend to seek niche products, while consumers with higher search costs gravitate toward standardized products. With two product designs, a reduction in search costs may raise prices for both if the endogenous shift in search modes---the composition effect---dominates the search-encouraging effect. Prices also respond asymmetrically to decreases in search costs, depending on where in the cost distribution the shock occurs.

\section*{Work in progress}

\subsection*{3. Too Much Unsold Stock?: Environmental Consequences of Price Commitment (with Kohei Kawaguchi and Kei Kawai)}
Resale price maintenance can be pro-competitive because it encourages sellers to carry inventory in the low-demand state. When realized demand falls short of that inventory, however, the seller cannot lower the price to clear the market, and unsold units are returned at an environmental cost. We build a dynamic oligopoly model with demand uncertainty, estimate it on monthly bookstore data, and ask how much unsold stock is due to price commitment and how a deposit-refund tax can reduce waste.

\subsection*{4. Pass-Through and Incidence with Product Returns (with Ce Bian and Liang Zhao)}
We study pass-through and tax incidence when products can be returned. Unlike the case without returns, a tax such as a platform commission or a sales tax can create negative pass-through and raise social welfare. When a higher price substantially raises the return probability, the seller reduces quantity supplied because the effective margin falls. This downward-sloping supply response can lower the equilibrium price after a tax increase, raising social surplus by reducing the return rate.
\end{document}
